\section{Concrete axioms for $A_\infty$ chiral algebras}
\label{sec:axioms-Ainfty-chiral}
\label{sec:ainfty_chiral_defn}

An $\Ainf$ chiral algebra is a $C_\ast(W(\SCchtop))$-algebra presented in coordinates: multilinear operations $\{m_k\}_{k \ge 1}$ satisfying sesquilinearity, the Stasheff identities (from Stokes on $\FM_k(\C)$), and symmetry.

\subsection{Underlying data and parameters}
\label{subsec:underlying-data}
\begin{definition}[$A_\infty$ chiral algebra]
\label{def:ainfty_chiral}
Fix a ground field $k$ of characteristic $0$. Let $(A,Q)$ be a cochain complex (cohomological grading), with $Q$ of degree $+1$.
We equip $A$ with:
\begin{itemize}
  \item a degree $0$ endomorphism $\partial\colon A\to A$ (holomorphic translation) with $[Q,\partial]=0$,
  \item a distinguished unit $\mathbf 1\in A^0$ with $Q\mathbf 1=0$ and $\partial \mathbf 1=0$,
  \item multilinear operations (for $k\ge 1$)
  \begin{equation}
    \label{eq:mk-type}
    m_k \colon A^{\otimes k}\longrightarrow A((\lambda_1))\cdots((\lambda_{k-1})),
    \qquad \deg m_k = 2-k,
  \end{equation}
  where $A((\lambda))$ denotes the space of formal Laurent series in $\lambda$ with coefficients in $A$, and $A((\lambda_1))\cdots((\lambda_{k-1}))$ is the iterated Laurent series ring.\footnote{We work in the completed topology so that infinite sums are allowed.}
\end{itemize}
We regard $\lambda_1,\dots,\lambda_{k-1}$ as spectral parameters associated to the \emph{differences} of the holomorphic coordinates of inputs $1,\dots,k-1$ relative to the $k$-th input (the ``rightmost'' slot).
The axioms imposed on this data are stated in Definition~\ref{def:sesquilinearity}--\ref{def:symmetry-grading} below.
\end{definition}

\begin{remark}[Categorical origin of the $A_\infty$ chiral algebra]
\label{rem:categorical-origin}
The $A_\infty$ chiral algebra $A$ defined above is the endomorphism algebra $A_b = \RHom_{\cC(J)}(b,b)$ of a vacuum object $b$ in the open-sector factorization category $\cC$ (Definition~\ref{def:oc-factorization-category}).  The multilinear operations $m_k$ arise from restricting the factorization functors of~$\cC$ to iterated self-$\Hom$'s of~$b$ (Theorem~\ref{thm:boundary-algebra-from-vacuum}).  The complete invariant of the open sector is the Morita class of~$\cC$, not the quasi-isomorphism class of~$A_b$: different compact generators $b, b'$ produce Morita-equivalent algebras $A_b \sim_M A_{b'}$ with equivalent module categories.
\end{remark}

\begin{definition}[Curved $A_\infty$ chiral algebra]
\label{def:curved-ainfty-chiral}
A \emph{curved $A_\infty$ chiral algebra} over a base $S$ (e.g., $S = \overline{\cM}_g$ for genus-$g$ modular theory) is the data of Definition~\ref{def:ainfty_chiral} together with:
\begin{itemize}
\item a \emph{curvature element} $m_0 \in \Gamma(S, \cO_S)$ of cohomological degree~$2$ (a section of the structure sheaf of the base, with no algebra inputs),
\end{itemize}
subject to the \emph{curved $A_\infty$ relations}: for each $n \geq 0$,
\[
\sum_{\substack{i+j = n+1 \\ i,j \geq 0}} \sum_{s=0}^{n-j} (-1)^{\epsilon(s,j)}\, m_i\bigl(\ldots, m_j(\ldots), \ldots\bigr) \;=\; 0,
\]
where the $j=0$ terms contribute $m_0$ insertions.  The first three cases are:
\begin{itemize}
\item \textbf{At $n=0$:}\; $m_1(m_0) = 0$.  The curvature element $m_0$ is closed under the differential~$m_1$.
\item \textbf{At $n=1$:}\; $m_1^2(x) = m_2(m_0, x) - m_2(x, m_0) = [m_0, x]_{m_2}$.  The differential does not square to zero; rather, it squares to the commutator with the curvature (the graded commutator reduces to the ordinary commutator because $|m_0| = 2$ is even).
\item \textbf{At $n=2$:}\; $m_1(m_2(x,y)) + m_2(m_1(x),y) + (-1)^{|x|}\,m_2(x,m_1(y)) + m_3(m_0,x,y) \pm m_3(x,m_0,y) \pm m_3(x,y,m_0) = 0$.  The binary product is $m_1$-linear up to homotopies controlled by $m_3$ and curvature insertions (setting $m_0=0$, this is the graded Leibniz rule: $m_1$ is a derivation of~$m_2$).
\end{itemize}
When $m_0 = \kappa(\cA) \cdot \omega_g$ for a scalar $\kappa$ and a base class $\omega_g$, the $n=1$ relation reduces to $d_{\barB}^2 = \kappa \cdot \omega_g$ (Volume~I, Theorem~D): the bar differential fails to square to zero by exactly the curvature obstruction, and this failure IS the non-formality at genus $g \geq 1$.  Setting $m_0 = 0$ recovers Definition~\ref{def:ainfty_chiral}.
\end{definition}

\begin{notation}[Spectral bookkeeping]
\label{not:spectral}
For $n\ge1$ and ordered inputs $a_1,\dots,a_n$, we always write $m_n(a_1,\dots,a_n)$ as a formal Laurent series in $(\lambda_1,\dots,\lambda_{n-1})$ following the convention above. If $I=[p+1,\dots,p+r]$ is a consecutive block, we denote
\[
\Lambda_I\;:=\;\lambda_{p+1}+\cdots+\lambda_{p+r-1},
\]
and for any index $j>p+r-1$ we define $\lambda^{(I)}_j:=\lambda_j$ (unchanged), whereas the \emph{effective parameter} for the collapsed block $I$ is $\Lambda_I$.
\end{notation}

\begin{remark}[Operadic formulation]
\label{rem:operadic-formulation}
\index{chiral endomorphism operad!operadic formulation}
The data of Definition~\ref{def:ainfty_chiral} admits a concise
operadic reformulation.  Define the \emph{chiral endomorphism operad}
$\cE\!nd^{\mathrm{ch}}_A$ as the nonsymmetric operad with components
\[
\cE\!nd^{\mathrm{ch}}_A(n)
:=
\Hom\bigl(A^{\otimes n},\,
A((\lambda_1))\cdots((\lambda_{n-1}))\bigr),
\]
and partial compositions $f \circ_i g$ defined by operadic insertion
with the block-substitution rule of Notation~\ref{not:spectral}:
the variable $\lambda_i$ in $f$ is replaced by the block sum
$\Lambda_I = \lambda_i + \cdots + \lambda_{i+m-2}$, and the
variables of $g$ are inserted consecutively.
Associativity of the composition follows from associativity of
consecutive block-substitution.  Then:
\begin{enumerate}[label=\textup{(\roman*)}]
\item An $A_\infty$ chiral algebra structure on $A$ is a degree-$+1$
  element $m = \sum_{k \ge 1} m_k \in \prod_{k \ge 1}
  \cE\!nd^{\mathrm{ch}}_A(k)$ satisfying $[m, m] = 0$, where
  the bracket is the Gerstenhaber bracket
  $[f,g] = f\{g\} - (-1)^{(|f|-1)(|g|-1)} g\{f\}$
  induced by operadic insertion.
\item The chiral Hochschild cochain complex
  $\mathrm{C}^\bullet_{\mathrm{ch}}(A, A)
  = \prod_{n \ge 0} \cE\!nd^{\mathrm{ch}}_A(n+1)[-n]$
  with differential $\delta = [m, -]$ is a brace dg algebra
  \textup{(}Volume~I, Theorem~\textup{\ref*{V1-thm:thqg-brace-dg-algebra}}\textup{)}.
\item The cohomology
  $\cZ^{\mathrm{der}}_{\mathrm{ch}}(A)
  := H^*(\mathrm{C}^\bullet_{\mathrm{ch}}(A, A), \delta)$
  is the \emph{chiral derived center} and carries a Gerstenhaber algebra
  structure.  It is the universal bulk of the holomorphic-topological
  theory with boundary~$A$
  \textup{(}Volume~I, Theorem~\textup{\ref*{V1-thm:thqg-swiss-cheese}}\textup{)}.
\end{enumerate}
This reformulation makes precise the connection between the
axiomatic definition and the operadic machinery of Volume~I,
namely the open/closed factorization construction developed in
Part~I.
\end{remark}

\subsection{Sesquilinearity and unitality axioms}
\label{subsec:sesqui}

\begin{definition}[Sesquilinearity and unitality]
\label{def:sesquilinearity}
The $m_k$ satisfy the following sesquilinearity relations (generalizing the vertex-algebra relations) for $k\ge 2$:
\begin{align}
\label{eq:sesqui-left}
m_k(a_1,\dots,\partial a_i,\dots,a_k)
&= -\,\lambda_i\cdot m_k(a_1,\dots,a_i,\dots,a_k), && 1\le i\le k-1,\\
\label{eq:sesqui-right}
m_k(a_1,\dots,a_{k-1},\partial a_k)
&= \Bigl(\lambda_1+\cdots+\lambda_{k-1}+\partial\Bigr)\, m_k(a_1,\dots,a_k).
\end{align}
\noindent
\emph{Convention.}\; In equations \eqref{eq:sesqui-left}--\eqref{eq:sesqui-right} we use the
\textbf{independent-parameter convention}: $\lambda_i$ is the spectral parameter
associated to the $i$-th input (the difference $z_i - z_k$ relative to the
rightmost slot), and each $\lambda_i$ is treated as an independent formal
variable. In this convention, sesquilinearity reads
$m_k(\ldots,\partial a_i,\ldots) = -\lambda_i\cdot m_k(\ldots)$ for
$1\le i\le k-1$. The alternative \emph{consecutive-difference convention}
used in Lemma~\ref{lem:sesquilinearity_explicit} defines $\lambda_i$ as the
difference $z_i - z_{i+1}$ between adjacent insertion points.
The two conventions are related by the change of variables
\begin{equation}\label{eq:indep-to-consec}
  \lambda_i^{\mathrm{indep}} = z_i - z_k
  = \sum_{j=i}^{k-1} (z_j - z_{j+1})
  = \sum_{j=i}^{k-1} \lambda_j^{\mathrm{consec}},
  \qquad 1\le i\le k-1.
\end{equation}
In the consecutive convention, the sesquilinearity identities
\eqref{eq:sesqui-left}--\eqref{eq:sesqui-right} become: for
$1\le i\le k-1$,
$m_k(\ldots,\partial a_i,\ldots) = -(\lambda_i+\cdots+\lambda_{k-1})
\cdot m_k(\ldots)$, and for $i=k$,
$m_k(\ldots,\partial a_k) = (\partial + \sum_{j=1}^{k-1} j\,\lambda_j)
\,m_k(\ldots)$.
(The factor $j\,\lambda_j$ arises from the change of variables: $\sum_{i=1}^{k-1}\lambda_i^{\mathrm{indep}} = \sum_{i=1}^{k-1}\sum_{j=i}^{k-1}\lambda_j^{\mathrm{consec}} = \sum_{j=1}^{k-1}\bigl(\sum_{i=1}^{j}1\bigr)\lambda_j^{\mathrm{consec}} = \sum_{j=1}^{k-1}j\,\lambda_j^{\mathrm{consec}}$, where the double sum is reordered by noting that $\lambda_j^{\mathrm{consec}}$ appears for each $i\le j$.)
See Lemma~\ref{lem:sesquilinearity_explicit} for the complete derivation.

For $k=1$, $m_1=Q$ and sesquilinearity is $[Q,\partial]=0$.
There is a (strict) unit $\mathbf 1$ such that
\begin{align}
\label{eq:unit-axioms}
m_1(\mathbf 1)&=0,\qquad \partial\mathbf 1=0,\\
\label{eq:unit-m2}
m_2(\mathbf 1,a)&=a,\qquad m_2(a,\mathbf 1)=(-1)^{|a|}a,\\
\label{eq:unit-higher}
m_k(\ldots,\mathbf 1,\ldots)&=0\quad \text{for all $k\neq 2$, whenever a unit is inserted}.
\end{align}
\noindent
\emph{Sign justification.}\; The sign $(-1)^{|a|}$ in
$m_2(a,\mathbf{1})=(-1)^{|a|}a$ arises from the bar desuspension: the
operations $m_k$ are defined as the components of a degree-$+1$
coderivation on the bar complex $T^c(s^{-1}\bar A)$, and the
desuspension $s^{-1}$ produces a Koszul sign $(-1)^{|a|}$ when
passing through the first input.  Concretely, the unshifted product
$\mu_2$ satisfies $\mu_2(a,\mathbf{1}) = a$, and the shifted operation
$m_2 = s \circ \mu_2 \circ (s^{-1} \otimes s^{-1})$ acquires the
sign $(-1)^{|a|}$ from the desuspension isomorphism acting on the
first tensor factor.
\end{definition}

\subsection{Planar $A_\infty$ relations with spectral substitution}
\label{subsec:ainfty-relations}

\begin{definition}[$A_\infty$ relations with spectral substitution]
\label{def:ainfty-relations}
For each $n\ge 1$, the following \emph{planar} $A_\infty$ identity holds (degree $3-n$):
\begin{multline}
\label{eq:ainfty-relation-raw}
\sum_{\substack{i+j=n+1\\ i,j\ge1}}\;\sum_{s=0}^{n-j}
(-1)^{\epsilon(s,j)}\;
m_i\!\Bigl(
a_1,\dots,a_s,\
m_j(a_{s+1},\dots,a_{s+j}\,;\,\lambda_{s+1},\dots,\lambda_{s+j-1})
\ \Big|\\
\Lambda_{[s+1,\dots,s+j]}
,\ a_{s+j+1},\dots,a_n;\ \bm\lambda^{\mathrm{out}} \Bigr)
\;=\;0.
\end{multline}
Here:
\begin{itemize}
  \item The sign $(-1)^{\epsilon(s,j)}$ is the Koszul sign for $A_\infty$ algebras in cohomological grading. The explicit formula is
  \[
    \epsilon(s,j) \;=\; |a_1|+\cdots+|a_s|,
  \]
  i.e., the parity of the total degree of the elements that the inner operation $m_j$ passes.%
  \footnote{Many references (Loday--Vallette, Keller) state the $A_\infty$
  identity on the \emph{desuspended} bar complex $s^{-1}\bar A$ with sign
  $(-1)^{rs+t}$ or equivalently
  $(-1)^{(j-1)(|\bar a_1|+\cdots+|\bar a_s|)}$ where
  $|\bar a_i|=|a_i|-1$. The present formula $\epsilon(s,j)=|a_1|+\cdots+|a_s|$
  (independent of~$j$) is the standard \emph{unsuspended} convention: it
  encodes the Koszul sign of passing a degree-$+1$ component of the bar
  differential past the first~$s$ elements. The two conventions produce
  the same identity $d_{\mathrm{bar}}^2=0$, differing only by signs absorbed
  into the definition of~$m_k$ via the desuspension isomorphism
  $m_k = \pm\, s\circ \mu_k\circ(s^{-1})^{\otimes k}$.
  Concretely, the two sign conventions yield the same
  $d_{\mathrm{bar}}^2 = 0$ because the desuspension isomorphism
  $s^{-1}$ converts the Koszul sign
  $(-1)^{(j-1)(|\bar a_1|+\cdots+|\bar a_s|)}$ (with $|\bar a_i| = |a_i|-1$)
  into $(-1)^{|a_1|+\cdots+|a_s|}$ after absorbing the factors of
  $(-1)^{j-1}$ from $s^{-1}$; the resulting quadratic identities are
  identical term by term.
  In particular, for $j=1$: $\epsilon(s,1)=|a_1|+\cdots+|a_s|$,
  so $Q=m_1$ is a graded derivation of every~$m_k$.}
  \item The notation $m_i(-\mid \Lambda_{[s+1,\dots,s+j]},-;\bm\lambda^{\mathrm{out}})$ means:
  the $j$ consecutive inputs $a_{s+1},\dots,a_{s+j}$ are first composed by $m_j$ with the internal spectral variables $(\lambda_{s+1},\dots,\lambda_{s+j-1})$ (relative to the $a_{s+j}$ slot), resulting in one effective input placed at position $s+1$ of the outer $m_i$ with \emph{effective spectral parameter} $\Lambda_{[s+1,\dots,s+j]}$. The other outer spectral variables are given by
  \[
    \bm\lambda^{\mathrm{out}} = \bigl(\lambda_1,\dots,\lambda_s,\ \Lambda_{[s+1,\dots,s+j]},\ \lambda_{s+j+1},\dots,\lambda_{n-1}\bigr),
  \]
  i.e.\ the inner block is replaced by its sum of parameters, while the rest is unchanged.
\end{itemize}
\begin{remark}[Spectral substitution principle]
\label{rmk:spectral-substitution}
Equation \eqref{eq:ainfty-relation-raw} encodes that \emph{when a block of inputs is fused first}, the correct spectral parameter inserted in the outer operation is the \emph{sum} of the inner $\lambda$'s, reflecting translation covariance and the convention that all $\lambda_i$ are differences relative to the last slot, the algebraic avatar of momentum conservation for formal distributions.
\end{remark}
\end{definition}

\subsection{Symmetry, grading, and parity}
\label{subsec:symmetry}

\begin{definition}[Symmetry and grading conventions]
\label{def:symmetry-grading}
We impose no symmetry requirement on the $m_k$ beyond the planar order (this is standard for $A_\infty$ structures). The Koszul rule governs signs in all axioms above. The grading convention is cohomological.
\end{definition}

\subsection{Regular/singular decomposition of $m_2$}
\label{subsec:reg-sing}
\label{def:Laurent-projection}
For $m_2(a,b)\in A((\lambda))$, we decompose
\[
m_2(a,b) = m_2^{\mathrm{reg}}(a,b)\;+\;m_2^{\mathrm{sing}}(a,b),
\]
where $m_2^{\mathrm{reg}}(a,b)\in A[\![\lambda]\!]$ and $m_2^{\mathrm{sing}}(a,b)\in \lambda^{-1}A[\![\lambda^{-1}]\!]$. We will use this in the cohomological PVA construction below.

\subsection{Consequences: product and bracket on cohomology}
\label{subsec:product-bracket}
\label{def:product}
\label{def:lambda-bracket}
As in Section~\ref{sec:Ainfty-to-PVA}, define for $[a],[b]\in H^\bullet(A,Q)$:
\begin{align}
\label{eq:prod-def}
[a]\cdot[b] &:= \Bigl[\, \frac12\bigl(m_2^{\mathrm{reg}}(a,b) + (-1)^{|a||b|}m_2^{\mathrm{reg}}(b,a)\bigr)\,\Bigr],\\
\label{eq:bracket-def}
\{[a]{}_\lambda [b]\} &:= \Bigl[\, \frac12\bigl(m_2^{\mathrm{sing}}(a,b) - (-1)^{(|a|+1)(|b|+1)}\; \mathrm{e}^{-(\lambda+\partial)\partial_\lambda}\,m_2^{\mathrm{sing}}(b,a)\bigr)\,\Bigr].
\end{align}
The exponential operator encodes the transformation $\lambda\mapsto -\lambda-\partial$ required by vertex skewsymmetry. (It acts on formal series by the standard rule $\mathrm{e}^{\alpha\partial_\lambda}f(\lambda)=f(\lambda+\alpha)$.)

\smallskip\noindent
\emph{Verification of the exponential operator.}\;
The operator $\mathrm{e}^{\alpha\partial_\lambda}$ acts on $f(\lambda)$ by
$f(\lambda+\alpha)$.  Setting $\alpha = -(\lambda+\partial)$ gives
$f(\lambda) \mapsto f(\lambda - \lambda - \partial) = f(-\partial)$,
which is precisely the vertex algebra skew-symmetry substitution
$\lambda \to -\lambda - \partial$.  The sign
$(-1)^{(|a|+1)(|b|+1)}$ is the Koszul sign for the $(-1)$-shifted
bracket: the shifted degrees are $|a|+1$ and $|b|+1$, and the bracket
has degree $-1$, so the transposition sign is
$(-1)^{(|a|+1)(|b|+1)}$.

\smallskip\noindent
\emph{Symmetrisation in the product.}\;
The symmetrisation in~\eqref{eq:prod-def} is necessary because $m_2$ is
not graded-commutative: it encodes the full OPE, not just the
commutative normally-ordered product.  The factor $\frac{1}{2}$ is
well-defined in characteristic~$0$ (a standing hypothesis on the ground
field, cf.\ Definition~\ref{def:ainfty_chiral}).  If
$m_2^{\mathrm{reg}}$ is already graded-commutative (as for the
Heisenberg algebra), the symmetrisation is redundant and
$[a]\cdot[b] = [m_2^{\mathrm{reg}}(a,b)]$.

We will show these are well-defined and satisfy the PVA axioms.

\subsection{Relationship to physical constructions}

\begin{lemma}[QME implies $A_\infty$ relations; \ClaimStatusProvedElsewhere]
\label{lem:QME_to_Ainfty}
For a physical realisation satisfying Theorem~\ref{thm:physics-bridge}:
The quantum master equation $\frac{1}{2}\{S_{\mathrm{eff}},S_{\mathrm{eff}}\}_{\mathrm{BV}} + \hbar\Delta S_{\mathrm{eff}}=0$, expanded in Feynman diagrams, is equivalent to the $A_\infty$ identities \eqref{eq:ainfty-relation-raw} for the operations $m_k$ extracted from $S_{\mathrm{eff}}$.
\end{lemma}

\begin{proof}
Expand $S_{\mathrm{eff}} = \sum_k \frac{1}{k!} S_k$ where $S_k$ encodes $k$-ary operations. The full QME $\tfrac{1}{2}\{S_{\mathrm{eff}},S_{\mathrm{eff}}\}_{\mathrm{BV}} + \hbar\Delta S_{\mathrm{eff}} = 0$ at each arity $n$ yields precisely the sum over compositions $m_i \circ m_j$ with $i+j=n+1$, matching \eqref{eq:ainfty-relation-raw}. The complete argument appears in \cite{CG17}*{\S5.4} and \cite{CDG20}*{\S3}.
\end{proof}

\begin{proposition}[Definition \ref{def:ainfty_chiral} captures the physical structure;
\ClaimStatusProvedHere]
\label{prop:physical_justification}
For physical realisations satisfying Theorem~\ref{thm:physics-bridge}, Definition~\ref{def:ainfty_chiral} formalizes the $A_\infty$ chiral algebra structure of Costello--Gwilliam--Gaiotto--Dimofte:

\begin{enumerate}[label=(\roman*)]
\item The operations $m_k$ encode the factorization product on $\Conf_k(\C)$, recovering the Costello--Gwilliam structure \cite{CG17}.

\item The $A_\infty$ identity \eqref{eq:ainfty-relation-raw} is the quantum master equation
$\frac{1}{2}\{S_{\mathrm{eff}}, S_{\mathrm{eff}}\}_{\mathrm{BV}} + \hbar \Delta S_{\mathrm{eff}} = 0$
expanded in Feynman diagrams.

\item The $\lambda$-bracket and form integrands match \cite{CDG20} and Khan--Zeng \cite{KZ25}: spectral parameters $\lambda_i$ encode derivatives via $e^{\lambda z}$ generating functions.

\item The definition recovers known structures:
$m_{k \geq 3} = 0$ gives ordinary PVAs (Example~\ref{ex:free_chiral});
one-loop-finite holomorphic theories give Gwilliam--Grady--Williams structures \cite{GGW21};
cohomology $H^\bullet(\A,Q)$ yields PVAs matching \cite{CDG20} (Theorem~\ref{thm:cohomology_PVA}).

\item The definition generalizes Tamarkin's chiral Hochschild complex \cite{Tam02}: it is the 3d analogue of 2d Hochschild theory from Poisson sigma models (Section~\ref{sec:chiral_hochschild_expanded}).
\end{enumerate}
\end{proposition}

\begin{remark}[Status]
\label{rem:status}
Costello--Dimofte--Gaiotto \cite{CDG20} state that ``an $A_\infty$ analog of a vertex algebra\ldots is a structure that has yet to be fully defined mathematically.''  Definition~\ref{def:ainfty_chiral} fills this gap; Proposition~\ref{prop:physical_justification} verifies it against the physical sources.
\end{remark}

\subsection{Factorization axiom and cluster limits}

As insertion points collide, the operations decompose according to the boundary stratification of $\FM_k(\C)$.  This \emph{cluster factorization} axiom is the bridge between the $m_k$ and the operadic structure of~$\OHT$.

\begin{theorem}[Cluster Factorization (Homotopy-Commutative Form);
\ClaimStatusProvedElsewhere]
\label{thm:cluster_factorization_homotopy}
Let $(\A, \{m_k\})$ be a logarithmic $\SCchtop$-algebra (a $C_\ast(W(\SCchtop))$-algebra whose weight forms factor as $\omega_k = \omega_k^{\mathrm{hol}} \otimes \omega_k^{\mathrm{top}}$ on $\FM_k(\C) \times \Conf_k(\R)$; the full definition is Definition~\ref{def:log-SC-algebra}) and let $\pi = \{S_1,\ldots,S_r\}$ be a partition of $\{1,\ldots,k\}$ with $|S_j| = k_j$. As points $\{z_i\}_{i \in S_j}$ collide to cluster centers $w_j$ (a limit within $\FM_k(\C)$), the operations satisfy
\begin{equation}
\label{eq:cluster_factorization}
\lim_{z_i \to w_j} m_k(O_1,\ldots,O_k) = \sum_{\sigma \in \Sh(\pi)} (-1)^{\varepsilon(\sigma;O)} \, m_r\big(O'_{\sigma(1)},\ldots,O'_{\sigma(r)}\big),
\end{equation}
where $O'_j := m_{k_j}(O_{i_1},\ldots,O_{i_{k_j}})$ for $S_j = \{i_1 < \cdots < i_{k_j}\}$, and spectral parameters compose via
\begin{equation}
\label{eq:spectral_composition}
\lambda_{\text{cluster},j} = \sum_{i \in S_j} \lambda_i + \partial_{w_j}.
\end{equation}
Here $\lambda_i$ for $i \in S_j$ denotes the inner spectral parameters relative to the cluster center; the sum runs over $|S_j| - 1$ independent parameters (the last being determined by the constraint $\sum_{i \in S_j} \lambda_i = \Lambda_j - \partial_{w_j}$).
The shuffle sum $\sum_{\sigma\in\Sh(\pi)}$ arises because the Swiss-cheese
operations carry a homotopy-commutative structure induced by the symmetric
group action on $\FM_k(\C)$. For the planar (non-commutative) specialization,
only the identity shuffle $\sigma=\mathrm{id}$ contributes, and the sum
reduces to a single term.
\end{theorem}

\begin{proof}[Proof outline; complete details in Appendix~\ref{app:FM_Stokes}]
The theorem follows from three ingredients:

\textbf{Step 1 (Geometric):} The boundary stratum $D_\pi \subset \partial \FM_k(\C)$ corresponding to partition $\pi$ is canonically isomorphic to
\[
D_\pi \cong \FM_r(\C) \times \prod_{j=1}^r \FM_{k_j}(\C),
\]
recording cluster centers $\{w_j\}$ and relative positions within each cluster.

\textbf{Step 2 (Analytic):} The form representative $\omega_k$ extends continuously to $D_\pi$, and its restriction $\omega_k|_{D_\pi}$ factors as the wedge product
\[
\omega_k|_{D_\pi} = \omega_r(O'_1,\ldots,O'_r) \wedge \bigwedge_{j=1}^r \omega_{k_j}(\{O_i\}_{i \in S_j}),
\]
up to shuffle signs.

\textbf{Step 3 (Spectral):} Holomorphic generating functions $e^{\sum_i \lambda_i z_i}$ degenerate to $e^{(\sum_{i \in S_j} \lambda_i) w_j}$ times relative terms, yielding (\ref{eq:spectral_composition}).

Complete details, including orientation/sign conventions, appear in Appendix \ref{app:FM_Stokes}.
\end{proof}

\begin{remark}[Nonabelian Hodge shadow; \ClaimStatusConjectured]
\label{rem:nonabelian_hodge_shadow}
The shadow tower of the chiral factor is expected to carry a nonabelian Hodge triple
$(\Theta_\cA, M_\cA, h_\cA)$:
the MC element $\Theta_\cA$ (flat side, $\bC$-direction), a shadow Higgs field $M_\cA$
(Higgs side), and a Shapovalov-type form $h_\cA$ (harmonic side), with the depth decomposition
$d = 1 + d_{\mathrm{arith}} + d_{\mathrm{alg}}$ arising as the Hodge
decomposition of this triple.  This structure is not used in the proof of the cluster factorization theorem; its precise formulation requires the nonabelian Hodge package developed in Volume~I.
\end{remark}

\paragraph{Remark 4.6 (Factorization as operadic composition).} Cluster factorization (\ref{eq:cluster_factorization}) is precisely the statement that $\A$ is an algebra over the chain operad $\OHT=C_\ast(W(\SCchtop))$: boundary faces in $\FM_k(\C)$ encode operadic compositions in the closed colour and ordered gluings along $\R$ encode open-colour compositions; Stokes' theorem (Section \ref{sec:FM_calculus}) supplies the coherence.

\subsection{Explicit coend formula and operadic composition}

The operadic definition of $\OHT$ has a concrete incarnation as a coend producing the operations~$m_k$.

\paragraph{Construction 4.7 (Coend formula for $m_k$).}
Using the $W$--model, the universal element determining $m_k$ is
\[
  m_k = \int^{T \in \mathrm{Trees}_k} \!\!\! C_\ast\!\bigl(W(\SCchtop)(T)\bigr)\otimes_{\mathrm{Aut}(T)}
  \bigotimes_{v \in \mathrm{Vert}(T)} \Bigl(
     C_\ast\!\bigl(\FM_{\mathrm{in}(v)}(\C)\bigr)\,\widehat{\otimes}\,
     C_\ast\!\bigl(\Conf_{\mathrm{out}(v)}(\R)\bigr)
  \Bigr)\, .
\]
Closed-colour factors record holomorphic collisions; open-colour factors record ordered
times; the $W$--parameters on internal edges control homotopies.
The coend exists because the source category $\mathrm{Trees}_k$ is finite
(finitely many labelled rooted trees with $k$ leaves) and all tensor
products are over a field of characteristic~$0$.
Remark~4.9 below clarifies: the vertex-wise factors in the formula are
determined by the output colour of each vertex, not by an independent
product decomposition.

\begin{example}[Binary Operation $m_2$ from the Coend]
For $k=2$, the only tree is the corolla with one internal vertex and two leaves. The coend formula yields
\[
m_2 = \cP^{\mathrm{ch}}(2) \otimes \cE_1(1),
\]
where $\cP^{\mathrm{ch}}(2)$ encodes the binary chiral operation (holomorphic collision of two points) and $\cE_1(1)$ provides the topological ordering (represented by a single edge/interval).

Explicitly, for operators $a, b \in \A$ at positions $(t_1, z_1), (t_2, z_2)$:
\[
m_2(a,b) = \int_{\Gamma \subset \FM_2(\C)} \int_{t_1 < t_2} K(z_1 - z_2, t_1 - t_2) \, a(t_1,z_1) \otimes b(t_2,z_2) \, dt_1 \, dt_2,
\]
where $K$ is the propagator kernel; for physical realizations it is supplied by Theorem~\ref{thm:physics-bridge}.
This integral formula applies to physical realizations satisfying Theorem~\ref{thm:physics-bridge}; for abstract logarithmic $\SCchtop$-algebras, the operation $m_2$ is defined directly by the operad action.
\end{example}

\paragraph{Remark 4.9 (Boardman--Vogt resolution details).}
Each internal edge $e$ of a tree $T$ carries a parameter $s_e\in[0,1]$ encoding the time of
composition. In the two--colour model one works with $W(\SCchtop)(T)$; there is no meaningful
vertex-wise factor $(P_{\mathrm{ch}}\rtimes E_1)$. Instead, vertex decorations are determined by the output
colour: closed vertices carry $\FM(\C)$--data; open vertices carry $\FM(\C)\times \Conf(\R)$--data. The
boundary identifications implement coherent sequential/simultaneous compositions.

\subsection{Sesquilinearity: Complete Proof}

\begin{lemma}[Sesquilinearity from Translation Equivariance;
\ClaimStatusProvedHere]
\label{lem:sesquilinearity_explicit}
Here we use the \textbf{consecutive-difference convention}: the spectral
parameter $\lambda_i$ is associated to the difference $z_i - z_{i+1}$
between adjacent insertion points. The translation $z_i\mapsto z_i+\varepsilon$
then shifts $\lambda_{i-1} = z_{i-1}-z_i$ by $-\varepsilon$ and
$\lambda_i = z_i - z_{i+1}$ by $+\varepsilon$.
This convention is related to the independent-parameter convention of
equations~\eqref{eq:sesqui-left}--\eqref{eq:sesqui-right} by the
change of variables~\eqref{eq:indep-to-consec}.

Let $\partial : \A \to \A$ denote holomorphic translation in $z$. Then the operations $m_k$ satisfy:
\begin{align}
\label{eq:sesqui_explicit_1}
m_k(\partial a_1, a_2, \ldots, a_k)
  &= -(\lambda_1 + \lambda_2 + \cdots + \lambda_{k-1})\, m_k(a_1, a_2, \ldots, a_k),\\
\label{eq:sesqui_explicit_i}
m_k(a_1, \ldots, a_{i-1}, \partial a_i, a_{i+1}, \ldots, a_k)
  &= -(\lambda_i + \lambda_{i+1} + \cdots + \lambda_{k-1})\, m_k(a_1, \ldots, a_k),
  \quad 2\le i\le k-1,\\
\label{eq:sesqui_explicit_k}
m_k(a_1, \ldots, a_{k-1}, \partial a_k)
  &= \Bigl(\partial + \sum_{j=1}^{k-1} j\,\lambda_j\Bigr)\, m_k(a_1, \ldots, a_k).
\end{align}
Note that \eqref{eq:sesqui_explicit_1} is the $i=1$ case of \eqref{eq:sesqui_explicit_i};
for $k=2$ it reduces to $-\lambda_1\,m_2$.
These are obtained from \eqref{eq:sesqui-left}--\eqref{eq:sesqui-right}
via the change of variables~\eqref{eq:indep-to-consec}.
\end{lemma}

\begin{proof}
The argument proceeds by analyzing the $\C$--translation
equivariance of the operadic structure maps of
$C_\ast(W(\SCchtop))$, which forces a precise relationship
between $\partial$ and the spectral parameters.

\medskip\noindent
\textbf{Step 1 (Translation action on configuration spaces).}
The group $(\C,+)$ acts on $\Conf_k(\C) = \{(z_1,\ldots,z_k) : z_i \neq z_j\}$ by
simultaneous translation: $\tau_\varepsilon(z_1,\ldots,z_k)
= (z_1+\varepsilon,\ldots,z_k+\varepsilon)$.  This action
descends to $\FM_k(\C)$ (which quotients by translations),
and is therefore \emph{trivial} on $\FM_k(\C)$.  Correspondingly,
the logarithmic weight forms $\omega_k \in
\Omega^\bullet_{\log}(\FM_k(\C))$ are translation-invariant.

Since $\FM_k(\C)$ is defined modulo translations, the operations
$m_k$ depend on the insertion points only through the
\emph{differences} $\lambda_i = z_i - z_{i+1}$ (consecutive-difference
convention).  Translating a single insertion
point $z_i \mapsto z_i + \varepsilon$ while keeping all others fixed is
\emph{not} the simultaneous translation (which is trivial), but
rather a non-trivial operation that changes the differences
$\lambda_{i-1} \mapsto \lambda_{i-1} - \varepsilon$ and
$\lambda_i \mapsto \lambda_i + \varepsilon$.

\medskip\noindent
\textbf{Step 2 (Sesquilinearity for the first input).}
Consider the operation
\[
  F(\varepsilon) := m_k(a_1(z_1 + \varepsilon), a_2(z_2),
  \ldots, a_k(z_k)).
\]
The left side is an element of $A((\lambda_1))\cdots((\lambda_{k-1}))$
depending on $\varepsilon$.  We compute $F(\varepsilon)$ in two
independent ways.

\emph{Via the input:}
$a_1(z_1+\varepsilon) = a_1(z_1) + \varepsilon\,\partial a_1(z_1)
+ O(\varepsilon^2)$, so
\begin{equation}\label{eq:sesqui-F-input}
  F(\varepsilon)
  = m_k(a_1,a_2,\ldots) + \varepsilon\, m_k(\partial a_1, a_2, \ldots)
    + O(\varepsilon^2).
\end{equation}

\emph{Via the spectral parameters:}
The shift $z_1 \mapsto z_1 + \varepsilon$ changes the first
spectral parameter $\lambda_1 = z_1 - z_2$ to $\lambda_1 +
\varepsilon$, while all other $\lambda_j$ ($j \geq 2$) are
unchanged.  Since $m_k$ is a formal Laurent series in
$\lambda_1$, we may expand directly:
\begin{equation}\label{eq:sesqui-F-spectral}
  F(\varepsilon)
  = m_k(a_1,\ldots)|_{\lambda_1 \mapsto \lambda_1 + \varepsilon}
  = m_k(a_1,\ldots)
    + \varepsilon\,\partial_{\lambda_1} m_k(a_1,\ldots)
    + O(\varepsilon^2).
\end{equation}
The substitution $\lambda_1 \mapsto \lambda_1 + \varepsilon$ is
a well-defined operation on formal Laurent series because it
preserves the order of the pole.
Comparing the $O(\varepsilon)$ terms of \eqref{eq:sesqui-F-input}
and \eqref{eq:sesqui-F-spectral}:
\[
  m_k(\partial a_1, a_2, \ldots, a_k)
  = \partial_{\lambda_1} m_k(a_1, \ldots, a_k).
\]
This derives the \emph{position-space} identity
$m_k(\partial a_1,\ldots) = \partial_{\lambda_1} m_k$
(the formal derivative with respect to $\lambda_1$).
The sesquilinearity identity \eqref{eq:sesqui_explicit_1} asserts
a different operation: multiplication by~$-\lambda_1$.  These two
are \emph{not} the same as operators on formal Laurent series.  The
correct relationship comes from the Borel correspondence.

The $A_\infty$ operations, when expressed in the divided-power generating function
$m_k = \sum a_{(n)}b\,\lambda_1^n/n!$ (the $\lambda$-bracket of vertex algebra
theory), satisfy the identity by mode calculation:
$(Ta)_{(n)} = -n\,a_{(n-1)}$ (translation covariance of modes) gives
\begin{align*}
  m_k(\partial a_1,\ldots;\lambda_1,\ldots)
  &= \textstyle\sum_{n\ge 0} (\partial a_1)_{(n)}\cdots\, \frac{\lambda_1^n}{n!}
  = \textstyle\sum_{n\ge 0} (-n)\, a_{(n-1)}\cdots\, \frac{\lambda_1^n}{n!}\\
  &= -\lambda_1 \textstyle\sum_{m\ge 0} a_{(m)}\cdots\, \frac{\lambda_1^m}{m!}
  = -\lambda_1\, m_k(a_1,\ldots;\lambda_1,\ldots).
\end{align*}
This establishes $m_k(\partial a_1,\ldots) = -\lambda_1\, m_k$ in the
consecutive convention for $k=2$ (where the two conventions coincide).
For $k\ge 3$, the mode argument yields the same identity with
$\lambda_1$ replaced by $\lambda_1^{\mathrm{indep}} = z_1 - z_k$, i.e.\
$m_k(\partial a_1,\ldots) = -\lambda_1^{\mathrm{indep}}\, m_k$.
Converting via \eqref{eq:indep-to-consec}:
$\lambda_1^{\mathrm{indep}} = \sum_{j=1}^{k-1}\lambda_j^{\mathrm{consec}}$,
giving \eqref{eq:sesqui_explicit_1}.

\medskip\noindent
\textbf{Step 3 (Sesquilinearity for the $i$-th input, $2\le i\le k$).}
The cleanest route uses the independent-parameter convention
$\lambda_j^{\mathrm{indep}} = z_j - z_k$ and then converts via
\eqref{eq:indep-to-consec}.

\emph{Case $2\le i\le k-1$.}\;
In the independent convention, translating $z_i\mapsto z_i+\varepsilon$
shifts only $\lambda_i^{\mathrm{indep}}$ (by $+\varepsilon$); the reference
point $z_k$ is unchanged, so there is no output translation.
Repeating the mode argument of Step~2 in the $i$-th variable gives
$m_k(\ldots,\partial a_i,\ldots) = -\lambda_i^{\mathrm{indep}}\,m_k$.
Converting:
$\lambda_i^{\mathrm{indep}} = \sum_{j=i}^{k-1}\lambda_j^{\mathrm{consec}}$,
which gives~\eqref{eq:sesqui_explicit_i}.

\emph{Case $i=k$.}\;
Translating $z_k\mapsto z_k+\varepsilon$ shifts every
$\lambda_j^{\mathrm{indep}} = z_j - z_k$ by $-\varepsilon$ and
also shifts the output position by $+\varepsilon$, contributing
$+\varepsilon\partial$.
The $O(\varepsilon)$ coefficient gives
$m_k(\ldots,\partial a_k)
 = \bigl(\partial + \sum_{j=1}^{k-1}\lambda_j^{\mathrm{indep}}\bigr)\,m_k$,
which is~\eqref{eq:sesqui-right} in the independent convention.
Converting via~\eqref{eq:indep-to-consec}:
\[
  \sum_{j=1}^{k-1}\lambda_j^{\mathrm{indep}}
  = \sum_{j=1}^{k-1}\sum_{\ell=j}^{k-1}\lambda_\ell^{\mathrm{consec}}
  = \sum_{\ell=1}^{k-1}\ell\,\lambda_\ell^{\mathrm{consec}},
\]
giving~\eqref{eq:sesqui_explicit_k}.

\medskip\noindent
\textbf{Step 4 (Convention reconciliation).}
Equations
\eqref{eq:sesqui_explicit_1}--\eqref{eq:sesqui_explicit_k} in the
consecutive convention and
\eqref{eq:sesqui-left}--\eqref{eq:sesqui-right} in the
independent convention are related by the linear change of variables
\eqref{eq:indep-to-consec}:
$\lambda_i^{\mathrm{indep}} = \sum_{j=i}^{k-1}\lambda_j^{\mathrm{consec}}$.
In particular:
for $1\le i\le k-1$, the independent-convention
identity $-\lambda_i^{\mathrm{indep}}\, m_k$ becomes
$-(\lambda_i+\cdots+\lambda_{k-1})^{\mathrm{consec}}\, m_k$;
and for $i=k$, the sum
$\sum_{j=1}^{k-1}\lambda_j^{\mathrm{indep}}$
equals $\sum_{\ell=1}^{k-1}\ell\,\lambda_\ell^{\mathrm{consec}}$
by interchanging the order of summation.
\end{proof}
\subsection{From $W(\mathsf{SC}^{\mathrm{ch,top}})$ to axioms and back}
\label{sec:operad-to-axioms}
\label{sec:equivalence}

\subsubsection{From the operad to explicit $m_k$ with spectral substitution}
\label{subsec:operad-to-mk}
Let $(A_{\mathsf{ch}},A_{\mathsf{top}})$ be a $C_\ast\!\bigl(W(\mathsf{SC}^{\mathrm{ch,top}})\bigr)$--algebra with underlying complex $A$ and translation operator $\partial$ (from the holomorphic \(\C\)--action). Choose, for each $k\ge1$, a fundamental chain $[\FM_k(\C)]\in C_\ast(\FM_k(\C))$ and similarly for $E_1$ (recording the $t$--ordering). Define
\[
m_k(a_1,\dots,a_k)
:= \bigl\langle\ \rho\bigl([\FM_k(\C)]\times [E_1(\text{ordering})]\bigr)\,,\, a_1\otimes\cdots\otimes a_k\ \bigr\rangle
\ \in\ A((\lambda_1))\cdots((\lambda_{k-1})),
\]
where $\rho$ is the operad action and the \(\lambda\)--dependence is the generating function of the log--forms (residue degrees) on $\FM_k(\C)$ prescribed by the operad.
Then:
\begin{enumerate}[label=(\alph*)]
  \item sesquilinearity \eqref{eq:sesqui-left}--\eqref{eq:sesqui-right} follows from equivariance under translations in $\C$ (the \(\C\)--action generates $\partial$ and gives \(\lambda\)-multiplications),
  \item the spectral substitution law in \eqref{eq:ainfty-relation-raw} is the image of the boundary identification of faces in $\partial\FM_k(\C)$ where a consecutive block collapses (the effective parameter is the sum of inner ones).
\end{enumerate}

\begin{proposition}[Operad $\Rightarrow$ axioms; \ClaimStatusProvedHere]
\label{prop:operad-implies-axioms}
The $m_k$ produced by the $W(\mathsf{SC}^{\mathrm{ch,top}})$--action satisfy the axioms in Section~\ref{sec:axioms-Ainfty-chiral}.
\end{proposition}

\begin{proof}
Let $(A_{\mathsf{ch}}, A_{\mathsf{top}})$ be a $C_\ast(W(\mathsf{SC}^{\mathrm{ch,top}}))$-algebra with operad action $\rho$, and let $m_k$ be defined by evaluation on fundamental chains as in \S\ref{subsec:operad-to-mk}:
\[
m_k(a_1,\ldots,a_k) := \bigl\langle \rho\bigl([\FM_k(\C)] \times [E_1(\mathrm{ord})]\bigr),\, a_1 \otimes \cdots \otimes a_k \bigr\rangle.
\]
We verify each axiom of Definition~\ref{def:ainfty_chiral}.

\medskip
\textbf{(I) Unitality.}
The operadic unit in $\mathsf{SC}^{\mathrm{ch,top}}$ is the identity element $\mathbf{1} \in \mathsf{SC}^{\mathrm{ch,top}}(1)$, represented by the point $\FM_1(\C) = \mathrm{pt}$. By the operadic unit axiom, $\rho(\mathbf{1}) = \mathrm{id}_A$, so $m_1(\mathbf{1}) = Q(\mathbf{1}) = 0$ (since the differential on the operad restricts to zero on the unit). For $m_2$: the operadic composition $\mathbf{1} \circ_1 c_2 = c_2 = c_2 \circ_2 \mathbf{1}$ (inserting the unit into either slot of the binary corolla) gives $m_2(\mathbf{1}, a) = a$ and $m_2(a, \mathbf{1}) = (-1)^{|a|} a$, where the sign arises from the Koszul rule for the degree-$(-1)$ operation $m_2$ passing element $a$. For $k \geq 3$: the operadic composition $c_k \circ_i \mathbf{1}$ factors through $c_{k-1}$ by the operadic unit law, but the corollas $c_k$ for distinct $k$ are independent generators of $W(\mathsf{SC}^{\mathrm{ch,top}})$, and the unit insertion collapses $c_k \circ_i \mathbf{1}$ to a face of the $W$-construction with Boardman--Vogt parameter $s = 0$, which is $\partial$-exact. Hence $m_k(\ldots, \mathbf{1}, \ldots) = 0$ for $k \geq 3$. This establishes~\eqref{eq:unit-axioms}--\eqref{eq:unit-higher}.

\medskip
\textbf{(II) Sesquilinearity from translation equivariance.}
The group $\C$ acts on $\FM_k(\C)$ by simultaneous translation $z_i \mapsto z_i + w$, and this action is compatible with the operad structure. The operator $\partial \colon A \to A$ is the infinitesimal generator of the $\C$-action on $A$ via $\rho$.

For the $i$-th input ($1 \leq i \leq k-1$), consider translating only the $i$-th insertion point: $z_i \mapsto z_i + \varepsilon$. This changes the relative coordinate $\lambda_i = z_i - z_{i+1}$ by $+\varepsilon$, while the generating function for $m_k$ transforms as $m_k \mapsto e^{-\varepsilon \lambda_i} m_k$ (the spectral parameter $\lambda_i$ is conjugate to $z_i - z_{i+1}$). Taking the infinitesimal ($\varepsilon \to 0$) and using $\partial a_i = \lim_{\varepsilon \to 0}(a_i(z_i+\varepsilon) - a_i(z_i))/\varepsilon$:
\[
m_k(a_1,\ldots,\partial a_i,\ldots,a_k) = -\lambda_i \cdot m_k(a_1,\ldots,a_k), \qquad 1 \leq i \leq k-1,
\]
establishing~\eqref{eq:sesqui-left}. For the last input ($i = k$): shifting $z_k \mapsto z_k + \varepsilon$ changes every $\lambda_j = z_j - z_{j+1}$ for $j < k$ (specifically, $\lambda_{k-1}$ shifts by $-\varepsilon$), and the output evaluation point shifts by $+\varepsilon$. The infinitesimal action on the output is $\partial$, while the total spectral shift is $e^{+\varepsilon(\lambda_1+\cdots+\lambda_{k-1})}$ (since each $z_j - z_k$ increases by $\varepsilon$). Combining:
\[
m_k(a_1,\ldots,\partial a_k) = (\lambda_1 + \cdots + \lambda_{k-1} + \partial)\, m_k(a_1,\ldots,a_k),
\]
establishing~\eqref{eq:sesqui-right}. This also follows from Lemma~\ref{lem:sesquilinearity_explicit} after reconciling the consecutive-difference and independent-parameter conventions.

\medskip
\textbf{(III) $A_\infty$ relations and spectral substitution from boundary faces.}
The differential on $C_\ast(W(\mathsf{SC}^{\mathrm{ch,top}}))$ is the boundary operator $\partial$ on chains of the Boardman--Vogt resolution. On the fundamental chain $[\FM_k(\C)]$, the boundary is $\partial[\FM_k(\C)] = \sum_S \epsilon_{D_S}\,[D_S]$, summing over collision divisors $D_S$ with $|S| \geq 2$. Since $\rho$ is a chain map ($\rho \circ \partial = d \circ \rho$, where $d$ on $\End(A)$ is the commutator with $Q$), applying $\rho$ to both sides gives:
\[
d\bigl(\rho([\FM_k(\C)])\bigr) = \sum_S \epsilon_{D_S}\, \rho([D_S]).
\]
Each face $D_S \cong \FM_{k-|S|+1}(\C) \times \FM_{|S|}^{\mathrm{red}}(\C)$ factorizes $\rho$ by operadic composition:
\[
\rho([D_S]) = \rho([\FM_{k-|S|+1}(\C)]) \circ_{\mathrm{ins}} \rho([\FM_{|S|}(\C)]) = m_{k-|S|+1}(\ldots, m_{|S|}(\ldots), \ldots).
\]
The identity $d^2 = 0$ on $\End(A)$ (equivalently $\partial^2 = 0$ on chains) then yields
\[
0 = d^2\bigl(\rho([\FM_k(\C)])\bigr) = \sum_{i+j=k+1}\sum_{s=0}^{k-j} (-1)^{\epsilon(s,j)}\, m_i(\ldots, m_j(\ldots), \ldots),
\]
which is the $A_\infty$ identity~\eqref{eq:ainfty-relation-raw}. (Only consecutive subsets $S$ contribute, by Theorem~\ref{thm:stokes_arnold}, Step~5.)

The \emph{spectral substitution} law follows from the geometry of the boundary identification. On $D_S$ for $S = \{s{+}1,\ldots,s{+}\ell\}$, the cluster center $u_S$ replaces the $\ell$ points of $S$ by a single effective point. The relative coordinate from $u_S$ to the next external point is $u_S - z_{s+\ell} = \sum_{p=s+1}^{s+\ell-1}(z_p - z_{p+1}) = \Lambda_S$, so the outer operation receives effective spectral parameter $\Lambda_S = \lambda_{s+1} + \cdots + \lambda_{s+\ell-1}$, while the remaining outer parameters are unchanged. This matches the spectral substitution in~\eqref{eq:ainfty-relation-raw} and Remark~\ref{rmk:spectral-substitution}.

\medskip
\textbf{(IV) Degree.}
The degree of $m_k$ is $2 - k$.  The bar complex $\barB(\cA) = T^c(s^{-1}\bar{\cA})$ carries a total differential $d_{\barB} = \sum_k d_k$ of cohomological degree $+1$.  Each $d_k$ is the coderivation with projection
$d_k|_{(s^{-1}\bar{\cA})^{\otimes k}} = s^{-1} \circ m_k \circ s^{\otimes k}\colon (s^{-1}\cA)^{\otimes k} \to s^{-1}\cA$.
The degree of this composition is $\deg(s^{-1}) + \deg(m_k) + k\cdot\deg(s) = -1 + \deg(m_k) + k$, and requiring this to equal $+1$ gives $\deg(m_k) = 2 - k$, in agreement with~\eqref{eq:mk-type}.  Concretely: $m_1$ (the BRST differential) has degree~$+1$, $m_2$ (the binary product) has degree~$0$, and $m_3$ (the first homotopy) has degree~$-1$.
\end{proof}

\subsubsection{Rectification: from axioms to a $W(\mathsf{SC}^{\mathrm{ch,top}})$--algebra}
\label{subsec:rectification}
Conversely, suppose $A$ carries $(Q,\partial,\{m_k\})$ satisfying the axioms of Section~\ref{sec:axioms-Ainfty-chiral}.

\begin{definition}[Tameness]\label{def:tameness}
We say that an $A_\infty$ chiral algebra $(A, \{m_k\})$ is \emph{tame} if:
\begin{enumerate}[label=(\alph*)]
\item the spectral parameter dependence of $m_k$ is polynomial in $\lambda_i$ (regular part) plus a Laurent tail in $\lambda_i^{-1}$ (singular part) with finitely many negative powers;
\item the operations $m_k$ are continuous with respect to the natural topology on $A((\lambda_1)) \cdots ((\lambda_{k-1}))$;
\item for sufficiently large $k$, the operations $m_k$ vanish (truncation).
\end{enumerate}
\end{definition}

\begin{remark}[Tameness for HT theories]
\label{rem:tameness-HT}
For the $A_\infty$ chiral algebras arising from holomorphic-topological
field theories (Theorem~\ref{thm:physics-bridge}), tameness is
automatic:
\begin{enumerate}[label=\textup{(\alph*)}]
\item The spectral parameter dependence comes from meromorphic
  expansions of propagator products, which are rational functions of the
  $\lambda_i$ with finitely many poles; this gives the
  polynomial$+$Laurent structure.
\item Continuity follows from the smooth dependence of the Feynman
  integrals on the test functions and the topology of the configuration
  spaces $\FM_k(\C) \times \Conf_k^{<}(\R)$.
\item Truncation ($m_k = 0$ for $k \gg 0$) holds when the loop order of
  each Feynman diagram contributing to $m_k$ is bounded by a function
  that eventually exceeds the loop budget.  Concretely, for a theory
  with $n$-valent interactions whose vertex valence strictly exceeds the
  $\mu$-slot count, the graph combinatorics force $m_k = 0$ for
  $k \gg 0$.  When the $\mu$-degree equals the vertex valence (as for
  Virasoro, where each cubic vertex has exactly 2 $\mu$-slots), no such
  bound applies: diagrams of arbitrarily high loop order contribute at
  every arity, the $A_\infty$ structure is infinite
  (cf.\ Proposition~\ref{prop:vir-truncation}), and tameness
  condition~(c) must be replaced by convergence of the bar codifferential
  on the completed cofree coalgebra.
\end{enumerate}
\end{remark}

\begin{theorem}[Axioms $\Rightarrow$ operad; \ClaimStatusProvedHere]
\label{thm:rectification}
Assume $A$ satisfies conditions~\textup{(a)--(b)} of Definition~\ref{def:tameness}. If additionally condition~\textup{(c)} holds (truncation), then there exists a (canonically defined up to contractible choice) $C_\ast\!\bigl(W(\mathsf{SC}^{\mathrm{ch,top}})\bigr)$-algebra structure on~$A$ inducing the given $m_k$ by evaluation on fundamental chains as above.  More generally, if condition~\textup{(c)} is replaced by convergence of the bar codifferential on the \emph{completed} cofree coalgebra $\widehat{T}^c(s^{-1}\bar{A})$ (as holds for the Virasoro and $W_N$ algebras), the same conclusion holds in the completed setting.
\end{theorem}

\begin{proof}
The proof proceeds in four steps: construct a coalgebra over the bar
cooperad $\mathbf{B}(\mathsf{SC}^{\mathrm{ch,top}})$ from the tame
$A_\infty$ data, apply the cobar functor to obtain a strict operad
algebra, verify agreement on fundamental chains, and invoke
homotopy--Koszulity for essential uniqueness.

\medskip
\textbf{Step 1 (Bar coalgebra from tame $A_\infty$ data).}
The tame $A_\infty$ chiral data $(Q, \partial, \{m_k\}_{k \ge 1})$
define a codifferential on the cofree conilpotent graded coalgebra
\[
  \mathcal{C}_A \;:=\; T^c\!\bigl(s^{-1}\bar{A}\bigr)
  \;=\; \bigoplus_{n \ge 1} (s^{-1}\bar{A})^{\otimes n},
\]
where $s^{-1}$ is desuspension and $\bar{A} = A / k\cdot\mathbf{1}$
is the augmentation ideal.  The codifferential
$d_{\mathcal{C}} \colon \mathcal{C}_A \to \mathcal{C}_A$
is the unique coderivation whose cogenerator projection
$\pi \circ d_{\mathcal{C}} \colon \mathcal{C}_A \to s^{-1}\bar{A}$
encodes the operations $m_k$: explicitly,
\[
  \pi \circ d_{\mathcal{C}}\bigl(
    [a_1 | \cdots | a_n]\bigr)
  \;=\; \sum_{k=1}^{n}\;\sum_{i=0}^{n-k}
    (-1)^{\epsilon(i,k)}\,
    s^{-1}\,m_k(a_{i+1}, \ldots, a_{i+k}).
\]
\emph{Why $d_{\mathcal{C}}^2 = 0$:} The $A_\infty$ identity
\eqref{eq:ainfty-relation-raw} is equivalent to the vanishing of
$\pi \circ d_{\mathcal{C}}^2$. Since $d_{\mathcal{C}}$ is a
coderivation on a cofree coalgebra, $d_{\mathcal{C}}^2$ is again a
coderivation; a coderivation on a cofree coalgebra is zero if and
only if its cogenerator projection is zero \cite{LV12}. Hence
$d_{\mathcal{C}}^2 = 0$.

\emph{Why $\mathcal{C}_A$ is a $\mathbf{B}(\mathsf{SC}^{\mathrm{ch,top}})$--coalgebra:}
The operations $m_k$ take values in iterated Laurent series
$A((\lambda_1))\cdots((\lambda_{k-1}))$, encoding the spectral
parameter dependence along $\FM_k(\C)$. The sesquilinearity axiom
\eqref{eq:sesqui-left}--\eqref{eq:sesqui-right} ensures that the
cogenerator maps intertwine the $\C$--translation action on
$\FM_k(\C)$ with $\partial$ on $A$.
The spectral substitution rule (Remark~\ref{rmk:spectral-substitution})
ensures that the cogenerator maps are compatible with the boundary
stratification of $\FM_k(\C)$: when a block $I$ collapses, the
effective spectral parameter is $\Lambda_I = \sum_{i \in I} \lambda_i$,
matching the face map of the bar cooperad.
Together, these equip $\mathcal{C}_A$ with the structure of a
conilpotent coalgebra over $\mathbf{B}(\mathsf{SC}^{\mathrm{ch,top}})$.
Tameness (Definition~\ref{def:tameness}) is essential here:
\begin{itemize}
  \item condition (a) (polynomial$+$Laurent spectral dependence)
    ensures that each cogenerator map lands in the algebraic
    completion $A((\lambda_1))\cdots((\lambda_{k-1}))$, so the
    pairing with chains on $\FM_k(\C)$ is well-defined;
  \item condition (b) (continuity) ensures the cogenerator maps
    extend to the completed tensor products that arise from
    compositions in the bar cooperad;
  \item condition (c) (eventual truncation: $m_k = 0$ for $k \gg 0$)
    ensures that the coalgebra is concentrated in finitely many
    arities, so the cofree coderivation is determined by finitely
    many cogenerator projections and all tree sums terminate.
\end{itemize}

\medskip
\textbf{Step 2 (Cobar functor produces a strict operad algebra).}
Apply the cobar functor $\Omega$ to $\mathcal{C}_A$:
\[
  \widetilde{A} \;:=\; \Omega(\mathcal{C}_A).
\]
By definition, $\widetilde{A}$ is the free
$\mathsf{SC}^{\mathrm{ch,top}}$--algebra on $s\,\mathcal{C}_A$
(the suspension of the cooperad coalgebra), equipped with the
differential $d_\Omega$ induced by both the internal differential
$d_{\mathcal{C}}$ and the cooperad cocompositions. Concretely,
$\widetilde{A}$ is an algebra over
$\Omega\mathbf{B}(\mathsf{SC}^{\mathrm{ch,top}})$.

Now invoke homotopy--Koszulity
(Theorem~\ref{thm:homotopy-Koszul}): the canonical counit
\[
  \epsilon \colon
  \Omega\mathbf{B}(\mathsf{SC}^{\mathrm{ch,top}})
  \;\xrightarrow{\;\sim\;}\;
  \mathsf{SC}^{\mathrm{ch,top}}
\]
is a quasi--isomorphism of two--colored dg operads.  Therefore
restriction of structure along $\epsilon$ converts the
$\Omega\mathbf{B}(\mathsf{SC}^{\mathrm{ch,top}})$--algebra
$\widetilde{A}$ into a
$C_\ast\!\bigl(W(\mathsf{SC}^{\mathrm{ch,top}})\bigr)$--algebra,
since $\Omega\mathbf{B}(\mathsf{SC}^{\mathrm{ch,top}})$
\emph{is} the Boardman--Vogt resolution
$W(\mathsf{SC}^{\mathrm{ch,top}})$ (the cobar of the bar is
the canonical cofibrant replacement of the operad \cite{BM03}).
Denote the resulting operad algebra structure also by
$\widetilde{A}$.

\medskip
\textbf{Step 3 (Agreement on fundamental chains).}
We verify that the $C_\ast\!\bigl(W(\mathsf{SC}^{\mathrm{ch,top}})\bigr)$--algebra
$\widetilde{A}$ recovers the original operations $m_k$ on
fundamental chains.  The cobar--bar counit
$\eta_A \colon A \to \Omega(\mathcal{C}_A) = \widetilde{A}$
embeds $A$ as the weight-$1$ component (single-vertex trees)
of the cobar complex. For the $k$--corolla $c_k$, the
structure map of $\widetilde{A}$ evaluated on the fundamental
chain $[\FM_k(\C)] \times [\mathrm{ord}]$ reduces, on
weight-$1$ elements, to the cogenerator projection of
$d_{\mathcal{C}}$:
\[
  \widetilde{m}_k(a_1, \ldots, a_k)
  \;=\; \pi \circ d_{\mathcal{C}}
    \bigl([a_1 | \cdots | a_k]\bigr)
  \;=\; m_k(a_1, \ldots, a_k).
\]
The spectral parameters are preserved because the
cogenerator maps encode exactly the Laurent coefficients of
the $\FM_k(\C)$--forms, and the bar--cobar adjunction respects
the topology on iterated Laurent series (which is where
tameness condition~(b) is used).

For composite trees $T$ with internal edges, the structure map
$\rho(T)$ of $\widetilde{A}$ is determined inductively by the
operad composition from the corolla data $\{m_k\}$. The
compatibility conditions are:
\begin{enumerate}[label=(\roman*)]
  \item The $A_\infty$ relations \eqref{eq:ainfty-relation-raw}
    ensure $d_{\mathcal{C}}^2 = 0$, which is exactly the
    statement that $\rho \circ \partial_{\FM} = d_A \circ \rho$
    on $C_\ast(\FM_k(\C))$ (boundary maps to compositions).
  \item Sesquilinearity \eqref{eq:sesqui-left}--\eqref{eq:sesqui-right}
    ensures that $\rho$ intertwines the $\C$--action on $\FM_k(\C)$
    with $\partial$ on $A$.
  \item The spectral substitution rule ensures that boundary
    faces $D_I \cong \FM_r(\C) \times \FM_{|I|}(\C)$ map
    correctly under $\rho$, with the block--sum
    $\Lambda_I = \sum_{i \in I} \lambda_i$ as effective parameter.
\end{enumerate}
These are precisely the axioms of a
$C_\ast(W(\mathsf{SC}^{\mathrm{ch,top}}))$--algebra
(cf.\ the forward direction, Proposition~\ref{prop:operad-implies-axioms}).

\medskip
\textbf{Step 4 (Uniqueness via Quillen equivalence).}
It remains to show that the
$C_\ast\!\bigl(W(\mathsf{SC}^{\mathrm{ch,top}})\bigr)$--algebra
structure is unique up to contractible ambiguity. By
Theorem~\ref{thm:homotopy-Koszul} and
Theorem~\ref{thm:bar-cobar-adjunction}, the bar--cobar adjunction
\[
  \Omega \colon
  \mathrm{Coalg}_{\mathbf{B}(\mathsf{SC}^{\mathrm{ch,top}})}
  \rightleftarrows
  \mathrm{Alg}_{\mathsf{SC}^{\mathrm{ch,top}}}
  :\! \mathbf{B}
\]
is a \emph{Quillen equivalence}: the unit
$A \to \Omega\mathbf{B}(A)$ and counit
$\mathbf{B}\Omega(C) \to C$ are weak equivalences.
Suppose $\rho, \rho'$ are two
$C_\ast(W(\mathsf{SC}^{\mathrm{ch,top}}))$--algebra structures on
$A$ that agree on fundamental chains (i.e., induce the same $m_k$).
Both structures produce the same bar coalgebra $\mathcal{C}_A$
(since the bar coalgebra is determined by the cogenerator
projections, which are the $m_k$). The Quillen equivalence ensures
that the cobar $\Omega(\mathcal{C}_A)$ is the unique (up to weak
equivalence) operad algebra recovering $\mathcal{C}_A$ via the bar
construction. More precisely, the mapping space
\[
  \mathrm{Map}_{\mathrm{Alg}}^{\mathbf{B}\text{-equiv}}
  \!\bigl(\widetilde{A}, \widetilde{A}\bigr)
\]
of self-equivalences of $\widetilde{A}$ that induce the identity on
$\mathbf{B}(\widetilde{A}) \cong \mathcal{C}_A$ is contractible:
the derived unit $\eta \colon \mathrm{id} \xrightarrow{\sim}
\Omega \circ \mathbf{B}$ of the Quillen equivalence provides a
canonical contraction.
(Over a field of characteristic zero, the model category
$\mathrm{Alg}_{\mathsf{SC}^{\mathrm{ch,top}}}$ is right proper
and the mapping spaces are computed by derived hom; contractibility
of the fiber follows from the fact that $\eta_A$ is a weak
equivalence \cite{Val07, BM03}.)

\medskip
\textbf{Equivalence of categories.}
The forgetful functor $U$ (extracting $m_k$ from
$\rho([\FM_k(\C)] \times [\mathrm{ord}])$) is the construction of
Proposition~\ref{prop:operad-implies-axioms}. The bar--cobar
reconstruction of Steps~1--3 provides an inverse $R$ with
$U \circ R \cong \mathrm{id}$ (by Step~3) and
$R \circ U \simeq \mathrm{id}$ (by Step~4, uniqueness up to
contractible ambiguity). Hence the homotopy categories are
equivalent:
\[
  \mathrm{Ho}\bigl(\text{tame axiomatic }A_\infty\text{ chiral algebras}\bigr)
  \;\simeq\;
  \mathrm{Ho}\bigl(W(\mathsf{SC}^{\mathrm{ch,top}})\text{-algebras with
  }\partial,\,\mathbf{1}\bigr).
  \qedhere
\]
\end{proof}

\begin{remark}[Lagrangian skeleton and Stokes origin of the $A_\infty$ relations]
\label{rem:lagrangian-skeleton}
The rectification theorem (Theorem~\ref{thm:rectification}) and
its converse (Proposition~\ref{prop:operad-implies-axioms}) have a
geometric interpretation through the Lagrangian correspondence
$\mathfrak{S}_b = \mathcal{L}\times_{\mathcal{M}}\mathcal{L}$
of Remark~\ref{rem:lagrangian-geometry}.

The $A_\infty$ relations \eqref{eq:ainfty-relation-raw} are
Stokes' theorem on the boundary strata of this correspondence.
Each codimension-$1$ face of $\FM_k(\C)$ is a collision stratum
of the Lagrangian (a locus where two or more sheets of
$\mathcal{L}$ meet transversally), and the sum of residues
along these faces vanishes by Stokes, which is the Stasheff
identity $\sum m_i\circ m_j = 0$.  The spectral substitution
law (Remark~\ref{rmk:spectral-substitution}) is the cluster
structure of the collision strata: when a block of points
collapses, the Lagrangian sheets merge and the effective
spectral parameter is the sum of the inner ones, reflecting
the additivity of the shifted symplectic pairing on strata.

Sesquilinearity \eqref{eq:sesqui-left}--\eqref{eq:sesqui-right}
is $\C$-equivariance of the Lagrangian embedding: the translation
action on $\FM_k(\C)$ preserves $\mathcal{L}$, and the spectral
parameters record the infinitesimal action on the normal bundle.
Tameness (Definition~\ref{def:tameness}) is the algebraicity
of the correspondence: the polynomial$+$Laurent structure
is finite-dimensionality of the fibers, and truncation is
properness of the projection.

In this language, rectification is the \emph{Lagrangian skeleton
theorem}: the self-intersection data of $\mathcal{L}$
determines the correspondence $\mathfrak{S}_b$ and hence the
$\mathsf{SC}^{\mathrm{ch,top}}$-algebra, just as the
Steinberg variety is determined by its moment-map fibers
and recovers the Hecke algebra from the convolution product
on its Borel--Moore homology.
\end{remark}

\begin{proposition}[Rectification is the Lagrangian skeleton theorem;
\ClaimStatusProvedHere]
\label{prop:rectification-lagrangian-skeleton}
The rectification theorem \textup{(}recovering an
$\mathsf{SC}^{\mathrm{ch,top}}$-algebra from $A_\infty$ axioms\textup{)}
is the Lagrangian skeleton theorem: the self-intersection data
$\mathfrak{S}_b = \mathcal{L}\times_{\mathcal{M}}^h\mathcal{L}$
with its groupoid structure determines the Lagrangian correspondence,
and conversely.  The $A_\infty$ axioms are the minimal Lagrangian
skeleton data.
\end{proposition}

\begin{proof}
By Theorem~\ref{thm:homotopy-Koszul}, the operad
$\mathsf{SC}^{\mathrm{ch,top}}$ is homotopy-Koszul.  Combining this
with the bar-cobar adjunction (Vol~I, Theorem~A, the twisting morphism representability theorem) gives a Quillen
equivalence
\[
  \Omega\barB \colon
  \mathsf{SC}^{\mathrm{ch,top}}\text{-}\mathrm{Alg}
  \;\rightleftarrows\;
  A_\infty\text{-}\mathrm{Alg}
  \;\colon\; \barB
\]
between $\mathsf{SC}^{\mathrm{ch,top}}$-algebras and $A_\infty$
algebras satisfying the axioms of Section~\ref{sec:axioms-Ainfty-chiral}.
In the Lagrangian language: the $A_\infty$ axioms encode the
Lagrangian $\mathcal{L} \hookrightarrow \mathcal{M}$; the bar
complex $\barB$ passes to the self-intersection groupoid
$\mathfrak{S}_b = \mathcal{L}\times_{\mathcal{M}}^h\mathcal{L}$;
and the cobar functor $\Omega$ recovers the
$\mathsf{SC}^{\mathrm{ch,top}}$-algebra from the groupoid.  The
Quillen equivalence states that the $A_\infty$ axioms and the
$\mathsf{SC}^{\mathrm{ch,top}}$-algebra structure determine each
other: $\barB$ passes from the SC-algebra to the $A_\infty$ data,
$\Omega$ recovers the SC-algebra from the bar coalgebra, and the
two operations are inverse on the homotopy category.  In the
Lagrangian interpretation, this is the skeleton theorem: the
self-intersection data (groupoid) determines the correspondence
(SC-algebra), and conversely, because the bar-cobar adjunction
provides the independent algebraic content without circularity.
\end{proof}

\subsubsection{Compatibility with the PVA construction}
\label{subsec:compat-PVA}
The proof of the PVA structure on cohomology in Section~\ref{sec:axioms-Ainfty-chiral} can be interpreted operadically as the statement that the associated graded (holomorphic filtration) of $W(\mathsf{SC}^{\mathrm{ch,top}})$ is the free product of BD chiral and $E_1$, so the cohomological operations descend from the closed color; the bracket is the classical limit (residue) of the operadic composites, and associativity is the regular part of the $m_2$ identity.
